\documentclass[aspectratio=169]{beamer}
\usepackage[utf8]{inputenc}
\usepackage[T1]{fontenc}
\usepackage{booktabs}
\usepackage{listings}
\usepackage{xcolor}
\usepackage{textcomp}

\usetheme{Madrid}
\usecolortheme{whale}
\setbeamertemplate{navigation symbols}{}
\setbeamertemplate{footline}[frame number]

\definecolor{codebg}{RGB}{245,245,245}
\lstdefinestyle{shell}{
  basicstyle=\ttfamily\footnotesize,
  backgroundcolor=\color{codebg},
  breaklines=true,
  columns=fullflexible,
  frame=single,
  rulecolor=\color{gray!40},
  showstringspaces=false,
}
\lstset{style=shell}

\title{FASER V10}
\subtitle{Install, Build, Run --- a consolidated reference}
\author{Andr\'e Rubbia}
\institute{IPA Rubbia Group, ETH Z\"urich}
\date{\today}

\begin{document}

\begin{frame}
  \titlepage
\end{frame}

\begin{frame}{Outline}
  \tableofcontents
\end{frame}

\section{Pipeline overview}

\begin{frame}{Pipeline overview}
  \begin{itemize}
    \item \textbf{ConvertGENIE} / \textbf{ConvertFASERMC}: convert a GENIE or
      official FASERMC neutrino-interaction sample into the common
      \texttt{TPOEvent} format
    \item \textbf{FASERG4} (\texttt{faserps}): runs that through Geant4,
      writes per-event \texttt{TcalEvent} truth files
    \item \textbf{Batch} (\texttt{batchreco.exe}): reconstructs truth files
      into \texttt{TPORecoEvent} files
    \item \textbf{EvDisplay} / analysis / TauSearch: consumes the
      reconstructed output
    \item \texttt{run\_regression\_tests.py}: wraps the whole chain end to
      end for an automated simulate~$\to$~reconstruct~$\to$~compare check
  \end{itemize}
\end{frame}

\section{Known issue: proton-fraction bug}

\begin{frame}{Known issue: FASERG4 scintillator proton fraction (V9
    $\to$ V10)}
  \footnotesize
  \begin{itemize}
    \item FASERG4's scintillator material definition (3DCal tracker
      layers) had the wrong free-proton (H) vs.\ carbon-12 (C) content ---
      fixed between production versions \textbf{V9} and \textbf{V10}
    \item \textbf{Root cause}: \texttt{G4Material::AddElement()} is
      overloaded --- one variant takes an atom count (\texttt{G4int}),
      another a mass fraction (\texttt{G4double}). Declaring the
      stoichiometry counters as \texttt{G4double} silently selected the
      mass-fraction overload, so atom counts were read as mass fractions
    \item \textbf{Symptom}: free-proton and carbon-12 targets interacted at
      almost the same rate in V9 ($\approx$1:1) instead of the physically
      expected $\approx$1:12 ($\propto$ mass number $A$); V10 correctly
      gives $\approx$12:1
    \item \textbf{Effect}: neutrino-induced CC/NC rates rise
      $\sim$15--18\% (a more isoscalar target supplies more neutrons for
      the neutrino-preferred $d$-quark coupling), antineutrino rates stay
      roughly flat; overall 3DCal fiducial total $+4.4\%$
    \item also shifts scintillator density ($1.03\to1.06\;\mathrm{g/cm^3}$),
      hadronic interaction length, and dE/dx ($Z/A$ change); Birks' constant
      has \emph{not} yet been re-checked against the new composition
    \item full analysis: \texttt{FASERflux/faserg4\_proton\_fraction\_bug.tex}
      (Fabio Cufino, Andr\'e Rubbia)
  \end{itemize}
\end{frame}

\section{Known issue: muon-spectrometer field consistency}

\begin{frame}{Known issue: muon-spectrometer field consistency}
  \footnotesize
  \begin{itemize}
    \item \textbf{Motivation}: the muon spectrometer's magnetic field was
      implemented \emph{twice} --- once in Geant4 simulation
      (\texttt{MuonMagneticField}) and once in GenFit reconstruction
      (\texttt{GenMagneticField}) --- two independent hand-written
      implementations kept in sync only by hand, with no test proving they
      agreed and no check that either matched the real detector geometry
    \item \textbf{Fix}: extracted one shared, framework-independent
      function, \texttt{FASER::ComputeMuonSpectrometerField()}; both G4 and
      GenFit now call it as thin unit-converting adapters, so the two field
      \emph{formulas} can no longer silently diverge ---
      \texttt{MuonSpectrometerFieldTest} (gtest) checks G4 and GenFit agree
    \item \textbf{Real bug caught by the test}: outside the field's
      modelled envelope, GenFit's plain ``Magnet'' branch fell through to a
      stray non-zero fallback ($0.001$\,kG in $Y$) where G4 returned an
      exact zero --- a genuine, if tiny, pre-existing sim/reco divergence
    \item \textbf{Second problem}: the slit position deciding where the
      field flips sign / cuts to zero was a bare hardcoded literal
      duplicated in \textbf{four} places, agreeing only by convention, not
      construction --- a future geometry change could silently
      desynchronize the field model from the geometry
    \item \textbf{Fix}: \texttt{MagnetGeometryProbe} samples the
      \emph{actual} \texttt{TGeoShape} loaded from the real GDML to find
      where the iron truly ends and the slit truly begins, so
      reconstruction derives the slit position live instead of trusting a
      literal; \texttt{MagnetGeometryProbeTest} unit-tests the probe
      itself, including a case that deliberately mismatches the geometry
      to prove a future drift would actually be caught
  \end{itemize}
\end{frame}

\section{Git workflow: sandboxed commit locks}

\begin{frame}[fragile]{Git workflow: sandboxed commit locks}
  \footnotesize
  \begin{itemize}
    \item working from a sandboxed shell connected to the Mac, a
      \texttt{git commit} can leave stale
      \texttt{.git/index.lock} / \texttt{.git/HEAD.lock} /
      \texttt{.git/objects/maintenance.lock} files behind --- the sandbox
      can create them but sometimes can't \texttt{unlink} them afterward
      (\texttt{Operation not permitted})
    \item symptom: the \emph{next} \texttt{git commit} fails outright with
\begin{lstlisting}[language=bash]
fatal: cannot lock ref 'HEAD': Unable to create
  '.../.git/HEAD.lock': File exists.
\end{lstlisting}
    \item previously required manually running \texttt{rm -f} on the lock
      files from a real Terminal window on the Mac before the next commit
      could proceed
    \item \textbf{fixed}: a one-time delete-permission grant for the
      connected repo folder now lets stale locks be cleared directly from
      the same session --- no manual Terminal step needed before any
      future commit
  \end{itemize}
\end{frame}

\section{Install from scratch}

\begin{frame}[fragile]{1. Install from scratch}
\begin{lstlisting}[language=bash]
git clone https://github.com/rubbiaa/FASER.git
cd FASER
source setup.sh
\end{lstlisting}
  \begin{itemize}
    \item \texttt{setup.sh} auto-detects the site (Andr\'e's Mac, the Ubuntu
      \texttt{ryzen01} box, or lxplus) and sources the matching
      \texttt{thisroot.sh} / \texttt{geant4.sh}
    \item hands off to \texttt{common\_setup.sh}: CLHEP/Rave/GenFit paths,
      \texttt{LD\_LIBRARY\_PATH}, \texttt{CMAKE\_PREFIX\_PATH}, \texttt{PATH},
      and a sanity check for missing/stale variables
    \item defines \texttt{fb} (``faser build''): reruns
      \texttt{cmake --build \$HOMEFASER/build -j} from anywhere
    \item unrecognized machine? it prints exactly which \texttt{elif} branch
      to add
    \item full from-scratch ROOT/Geant4 install: \texttt{docs/INSTALL.md},
      \texttt{docs/GEANT4\_INSTALL.md}
  \end{itemize}
\end{frame}

\section{Build}

\begin{frame}[fragile]{2. Build}
\begin{lstlisting}[language=bash]
cmake -S . -B build -DCMAKE_BUILD_TYPE=RelWithDebInfo
cmake --build build -j
\end{lstlisting}
  \begin{itemize}
    \item first build also compiles CLHEP, Rave, GenFit, googletest, and
      Pythia8 (CMake superbuild) --- slow once, incremental afterward
    \item after \texttt{source setup.sh}, \texttt{fb} reruns the build step
      from anywhere
    \item executables land in \texttt{build/bin/}: \texttt{faserps},
      \texttt{batchreco.exe}, \texttt{batchreco\_detresp.exe},
      \texttt{dumphits.exe}, \texttt{evDisplay.exe},
      \texttt{ConvertGENIE.exe}, \texttt{Convert.exe},
      \texttt{CombineFluxes.exe}, the gtest binaries, \ldots
  \end{itemize}
\end{frame}

\begin{frame}{CMake build options}
  \centering
  \scriptsize
  \begin{tabular}{@{}p{3.6cm}p{1.1cm}p{6.3cm}@{}}
    \toprule
    \textbf{Option} & \textbf{Default} & \textbf{Meaning} \\
    \midrule
    \texttt{FASER\_BUILD\_EXTERNALS} & ON & Build CLHEP/Rave/GenFit/
      googletest/Pythia8 from source \\
    \texttt{FASER\_BUILD\_GEANT4\_PACKAGES} & ON & Build FASERG4 and
      FASERCalProtoG4 (needs Geant4) \\
    \texttt{FASER\_BUILD\_DISPLAY} & ON & Build the GenFit-based event
      display / muon spectrometer package \\
    \texttt{FASER\_BUILD\_TAUSEARCH} & ON & Build TauSearch \\
    \texttt{FASER\_BUILD\_TESTS} & ON & Build FASER's own gtest suite,
      registered with CTest \\
    \bottomrule
  \end{tabular}
  \par\medskip
  \footnotesize Full list and reuse-an-existing-dependency examples:
  \texttt{docs/INSTALL.md}
\end{frame}

\section{Where the data lives}

\begin{frame}{3. \texttt{\$FASERDATA}}
  \begin{itemize}
    \item every FASERG4/Batch executable resolves its data location through
      \texttt{\$FASERDATA} (\texttt{FASER::GetDataDir()}, in
      \texttt{CoreUtils/FaserDataDir.hh}) --- not a hardcoded relative path
      or an inter-directory symlink
    \item \texttt{GetDataDir()} creates a requested subdirectory on demand
      if it doesn't already exist
    \item \texttt{common\_setup.sh} defaults \texttt{FASERDATA} to
      \texttt{\$HOMEFASER/data} (gitignored) unless a per-site branch in
      \texttt{setup.sh} already exported it elsewhere (scratch, EOS, a
      dedicated checkout, \ldots)
  \end{itemize}
\end{frame}

\begin{frame}{\texttt{\$FASERDATA} subdirectories}
  \centering
  \scriptsize
  \begin{tabular}{@{}p{2.3cm}p{4.9cm}p{4.0cm}@{}}
    \toprule
    \textbf{Subdirectory} & \textbf{Contents} & \textbf{Written by / Read by} \\
    \midrule
    \texttt{faserG4/} & Per-event \texttt{TcalEvent} truth files &
      \texttt{faserps} / \texttt{batchreco.exe}, \texttt{dumphits.exe},
      \texttt{evDisplay.exe} \\
    \texttt{batch/} & Reconstructed \texttt{TPORecoEvent} files &
      \texttt{batchreco.exe} / analysis code, \texttt{summarize\_output.py} \\
    \texttt{GENIE/} & GENIE input samples (too large for git) &
      \texttt{fetch\_data.py} / \texttt{faserps} (\texttt{-{}-input-file}) \\
    \texttt{GDML/} & Detector geometry, exported every run &
      \texttt{faserps} / \texttt{run\_batchreco.py}'s default
      \texttt{-{}-geometry-file} \\
    \bottomrule
  \end{tabular}
\end{frame}

\section{Python wrapper scripts}

\begin{frame}{4. \texttt{run\_faserps.py} --- runs the Geant4 simulation}
  \begin{itemize}
    \item builds the V10-geometry macro in memory and pipes it straight to
      \texttt{faserps}' stdin --- no \texttt{.mac} file to keep in sync
    \item full prose docs (custom geometry, interactive UI, MuonDIS
      physics): \texttt{run\_faserps.md}
    \item default input sample is fetched automatically from a public
      CERNBox link the first time it's needed (sha256-verified)
  \end{itemize}
\end{frame}

\begin{frame}{\texttt{run\_faserps.py} flags (1/2)}
  \centering
  \tiny
  \begin{tabular}{@{}p{3.1cm}p{2.1cm}p{5.8cm}@{}}
    \toprule
    \textbf{Flag} & \textbf{Default} & \textbf{Meaning} \\
    \midrule
    \texttt{-{}-build-dir} & \texttt{build/} & CMake build dir containing
      \texttt{bin/faserps} \\
    \texttt{-{}-vis} & off & Launch the interactive Geant4 UI (no macro used) \\
    \texttt{-{}-print-macro} & off & Print the resolved macro text, exit \\
    \texttt{-{}-dry-run} & off & Print the command, don't execute \\
    \texttt{-{}-input-file} & default GENIE sample & ROOT input file; a
      relative custom value resolves against \texttt{FASERG4/} \\
    \texttt{-{}-start-event} & 0 & First event index to read \\
    \texttt{-{}-n-events} & 100 & Number of events (\texttt{/run/beamOn}) \\
    \texttt{-{}-muons} & off & Generate fixed-momentum muons instead of
      reading \texttt{-{}-input-file} \\
    \texttt{-{}-muon-momentum-gev} & 100.0 & Muon momentum [GeV] \\
    \bottomrule
  \end{tabular}
\end{frame}

\begin{frame}{\texttt{run\_faserps.py} flags (2/2)}
  \centering
  \tiny
  \begin{tabular}{@{}p{3.1cm}p{2.1cm}p{5.8cm}@{}}
    \toprule
    \textbf{Flag} & \textbf{Default} & \textbf{Meaning} \\
    \midrule
    \texttt{-{}-muondis} & off & Enable on-the-fly Pythia8 DIS per muon
      interaction (implies \texttt{-{}-muons}) \\
    \texttt{-{}-muondis-cross-section-bias} & 150.0 & Cross-section bias
      (\texttt{1} = unbiased) \\
    \texttt{-{}-muondis-q2min} & 1.0 & Minimum $Q^2$ [GeV$^2$] \\
    \texttt{-{}-muondis-interaction-log} & disabled & CSV path logging each
      interaction \\
    \texttt{-{}-muondis-pdf-set} & built-in proton PDF & Path under
      \texttt{FASERG4/input/} \\
    \texttt{-{}-muondis-xbjmin} & 0.0 & Minimum Bjorken-$x$ cut \\
    \texttt{-{}-muondis-debug} & off & Verbose MuonDIS debug output \\
    \texttt{-{}-tilt-deg} & $-4.5$ & Detector tilt about Y [deg] \\
    \texttt{-{}-shift-x-cm} / \texttt{-{}-shift-y-cm} & 45.0 / 24.0 & LOS shift
      [cm] \\
    \bottomrule
  \end{tabular}
\end{frame}

\begin{frame}{\texttt{run\_batchreco.py} --- runs the reconstruction}
  \begin{itemize}
    \item wraps \texttt{Batch/batchreco.exe} with named, validated flags
      instead of positional arguments
    \item full prose docs: \texttt{run\_batchreco.md}
  \end{itemize}
  \vspace{0.3cm}
  \centering
  \tiny
  \begin{tabular}{@{}p{3.1cm}p{3.0cm}p{4.9cm}@{}}
    \toprule
    \textbf{Flag} & \textbf{Default} & \textbf{Meaning} \\
    \midrule
    \texttt{-{}-run} & \emph{required} & Run number --- deliberately no
      default \\
    \texttt{-{}-min-event} / \texttt{-{}-max-event} & 0 / 99999999 & Event
      index range $[\mathrm{min},\mathrm{max})$ \\
    \texttt{-{}-mask} & none (all events) & \texttt{nueCC}, \texttt{numuCC},
      \texttt{nutauCC}, \texttt{nuNC}, or \texttt{nuES} \\
    \texttt{-{}-multi-thread} & off & Pass \texttt{-mt} to
      \texttt{batchreco.exe} \\
    \texttt{-{}-geometry-file} & \texttt{\$FASERDATA/\allowbreak
      GDML/\allowbreak FASERCAL\_V10.gdml} & GDML geometry to load \\
    \texttt{-{}-build-dir} & \texttt{build/} & CMake build dir containing
      \texttt{bin/batchreco.exe} \\
    \texttt{-{}-split} & 1 & Split the range into N parallel background jobs \\
    \texttt{-{}-dry-run} & off & Print the command(s), don't execute \\
    \bottomrule
  \end{tabular}
\end{frame}

\begin{frame}[fragile]{\texttt{run\_regression\_tests.py} --- simulate
    $\to$ reconstruct $\to$ compare}
  \begin{itemize}
    \item runs the real pipeline end to end for several cases
      (\texttt{nueCC}/\texttt{numuCC}/\texttt{nutauCC}/\texttt{nuNC}, plus
      \texttt{-{}-muons} and \texttt{-{}-muondis})
    \item compares aggregate output statistics against a committed golden
      JSON baseline
  \end{itemize}
\begin{lstlisting}[language=bash]
python3 run_regression_tests.py --record   # first time / deliberate change
python3 run_regression_tests.py            # every other time: compare
\end{lstlisting}
\end{frame}

\begin{frame}{\texttt{run\_regression\_tests.py} flags}
  \centering
  \scriptsize
  \begin{tabular}{@{}p{2.6cm}p{1.6cm}p{6.6cm}@{}}
    \toprule
    \textbf{Flag} & \textbf{Default} & \textbf{Meaning} \\
    \midrule
    \texttt{-{}-case} & all cases & Restrict to a golden case name
      (repeatable); see \texttt{-{}-list} \\
    \texttt{-{}-record} & off & Write the golden file(s) instead of comparing \\
    \texttt{-{}-rel-tol} & $10^{-9}$ & Relative tolerance for floating-point
      comparisons \\
    \texttt{-{}-build-dir} & \texttt{build/} & CMake build directory \\
    \texttt{-{}-python} & current interpreter & Interpreter used for the
      subprocess calls (point at a PyROOT-enabled one if needed) \\
    \texttt{-{}-list} & --- & List available case names, exit \\
    \bottomrule
  \end{tabular}
\end{frame}

\begin{frame}{\texttt{fetch\_data.py} \quad \& \quad
    \texttt{summarize\_output.py}}
  \scriptsize
  \textbf{\texttt{fetch\_data.py}} --- public-CERNBox-link downloads into
  \texttt{\$FASERDATA}, sha256-verified, used automatically by
  \texttt{run\_faserps.py}.

  \vspace{0.45cm}
  \centering
  \begin{tabular}{@{}p{2.3cm}p{1.6cm}p{6.6cm}@{}}
    \toprule
    \textbf{Flag} & \textbf{Default} & \textbf{Meaning} \\
    \midrule
    \texttt{-{}-force} & off & Re-download even if already present \\
    \texttt{-{}-dry-run} & off & Show what would be fetched, don't fetch \\
    \texttt{-{}-list} & --- & List manifest entries, exit \\
    \bottomrule
  \end{tabular}

  \vspace{0.4cm}
  \raggedright
  \textbf{\texttt{Tests/regression/summarize\_output.py}} --- PyROOT
  summarizer used internally by \texttt{run\_regression\_tests.py}
  (requires \texttt{-{}-truth-dir} and/or \texttt{-{}-reco-file}).

  \vspace{0.45cm}
  \centering
  \begin{tabular}{@{}p{2.3cm}p{2.6cm}p{5.6cm}@{}}
    \toprule
    \textbf{Flag} & \textbf{Default} & \textbf{Meaning} \\
    \midrule
    \texttt{-{}-dict-path} & CoreUtils' built dict & Path to the built ROOT
      dictionary \\
    \texttt{-{}-truth-dir} & none (skip) & Dir of truth files; needs
      \texttt{-{}-run}, \texttt{-{}-n-events} \\
    \texttt{-{}-reco-file} & none (skip) & Path to one
      \texttt{Batch-TPORecevent\_*.root} file \\
    \bottomrule
  \end{tabular}
\end{frame}

\section{Running it}

\begin{frame}[fragile]{5. A typical end-to-end run}
\begin{lstlisting}[language=bash]
source setup.sh
cmake -S . -B build -DCMAKE_BUILD_TYPE=RelWithDebInfo && cmake --build build -j

python3 run_faserps.py --n-events 500               # simulate (auto-fetch)
python3 run_batchreco.py --run 10000 --mask numuCC   # reconstruct numuCC
build/bin/evDisplay.exe 10000 numuCC                 # look at the result
\end{lstlisting}
  or, to exercise the whole pipeline and check it against the committed
  baseline:
\begin{lstlisting}[language=bash]
python3 run_regression_tests.py
\end{lstlisting}
\end{frame}

\begin{frame}[fragile]{6. Tests}
\begin{lstlisting}[language=bash]
ctest --test-dir build --output-on-failure
\end{lstlisting}
  \begin{itemize}
    \item FASER's own gtest suite (\texttt{MuonSpectrometerFieldTest},
      \texttt{MagnetGeometryProbeTest}, \ldots) builds automatically
      whenever \texttt{FASER\_BUILD\_TESTS} and
      \texttt{FASER\_BUILD\_GEANT4\_PACKAGES} are both \texttt{ON} (the
      default)
    \item the heavier simulate-reconstruct-compare regression suite is
      \texttt{run\_regression\_tests.py}, described above and in
      \texttt{docs/REGRESSION\_TESTS.md}
  \end{itemize}
\end{frame}

\section{Further reading}

\begin{frame}{Further reading}
  \begin{itemize}
    \item \texttt{docs/INSTALL.md} --- the full CMake build: prerequisites,
      every option, troubleshooting a clean build
    \item \texttt{docs/GEANT4\_INSTALL.md} --- building/installing Geant4
      itself
    \item \texttt{run\_faserps.md} / \texttt{run\_batchreco.md} --- prose
      documentation for those two wrappers
    \item \texttt{docs/README\_MuonDIS.md} --- MuonDIS physics and every
      \texttt{/physics/muondis/...} option
    \item \texttt{docs/REGRESSION\_TESTS.md} --- the regression suite's
      full design, what's verified vs.\ not, and open items
  \end{itemize}
\end{frame}

\end{document}
